\itemhead{复习题一 · 第10题（教材p.30）}{XJ-B1-C01-RV-P030-Y0307}
\label{XJ-B1-C01-RV-P030-Y0307}
{\footnotesize 湘教版教材p.30 / PDF第37页 / 复习题一 / 第10题}
\field{原题}{用全称量词或存在量词的符号表述命题：“任意三角形 $ABC$ 都有外接圆。”}
\field{参考答案}{$\forall\triangle ABC,\exists$ 圆 $O$，使 $A,B,C$ 都在圆 $O$ 上。存在的圆可随三角形而改变，不能交换两个量词次序。}
{\small \field{关联知识}{L03 全称、存在量词及真假的判定（教材帮p.38、39 / PDF47、48）}
\field{方法与易错点}{先确认量词、作用对象及论域，再采用任意性证明或实例反例。 易错：用几个满足的数证明全称命题；把“存在”理解为“唯一存在”。}
\field{答案依据}{答案册相关题组：PDF第8、9页（册内页7、8）；按当前题库原题逐项重算，不按旧版题号机械对套。}}
